\subsection{Two characteristic-invariant cuts}
\label{subsec:ha-two-hinge}

The seven-cut threshold above concerns ten fixed radial profiles. A
different obstruction applies to any two affine hinges that are
constant along their own facet characteristic directions. Affine
offsets, dependence on the facet-edge coordinate, and unequal cut
speeds are allowed. For a nonnegative function $f$ on the disjoint
labeled facet union, put
\[
 R(f)=\sup_\gamma\int f\,d\mu_\gamma,
\]
where the supremum runs over all finite maximizing words, equivalently
their weak occupation limits.

\begin{proposition}[Two invariant hinges]\label{prop:ha-two-hinge}
On one HKO facet $F_i$, let $f_i(x)=(D\cdot x)_+$, with $f$ zero on the
other labeled facets and $D\cdot\mathbb{J}a_i=0$. Then
\begin{equation}\label{eq:ha-amplification}
 R(f)\geq2\nu(f).
\end{equation}
Consequently, retain all ten original inequalities of $H$ and add at
most two rows $(a_i+tD)\cdot x\leq1$ satisfying this characteristic
invariance. Their one-sided first logarithmic systolic variation at
$t=0$ is nonpositive, whether the rows concern the same facet or
different facets.
\end{proposition}

Appendix~\ref{sec:ha-two-hinge-proof} proves the one-hinge bound from
four literal shortest-characteristic families and a scalar hinge
inequality, then handles both placements of two cuts. This is a
first-order obstruction for invariant affine hinges. It does not prove
local maximality among twelve-facet bodies, decide zero-derivative
directions at higher order, or cover non-invariant hinges.
